\section{Problem and Experimental Setting}
\label{sec:setting}

\paragraph{Streams.}
Inputs are $x\sim\mathcal N(0,I_{16})$. We use two conditions, both fixed before running.
In \emph{random-teacher} ($C{=}4$ classes), task $k$ labels $x$ by $\arg\max_c (W_k x)_c$ with an
independent Gaussian teacher $W_k$. In \emph{label-permutation} ($C{=}10$), a single teacher is shared
and task $k$ applies a random label permutation. We use 10 classes because with 4 classes only 24
permutations exist, so a 50-task stream would repeat them and the probe task would not be new. A runtime
check verifies that the probe permutation never occurs in training. Each stream has 50 training tasks of
200 steps with batch size 16, and one additional task that is never trained in the stream and serves as
the probe.

\paragraph{Model and optimizers.}
The model is a one-hidden-layer MLP (16--32--$C$) with BatchNorm (momentum 0.1), ReLU and dropout 0.1,
trained with mean cross-entropy and no replay, regulariser or task identity. The base optimizer is Adam
($\beta_1{=}0.9$, $\beta_2{=}0.999$, $\eps{=}10^{-8}$, no weight decay, no AMSGrad). Its learning rate is
selected on two development seeds by online training accuracy over the whole stream, from
$\{10^{-3},3{\cdot}10^{-3},10^{-2},3{\cdot}10^{-2}\}$. A single one-sided extension was allowed if the
choice fell on a grid edge. The selected values were \RTTuneAdamLr{} (random teacher) and \LPTuneAdamLr{} (label permutation), both
interior grid points. Adam with $\beta_1{=}\beta_2{=}0.9$ and SGD with momentum 0.9 were tuned in the same way and are
reported as secondary families (Appendix~\ref{app:tuning}).

\paragraph{Probe protocol.}
From each starting point we train for 200 steps on the probe task and evaluate held-out probe accuracy
every 10 steps. The three starting points are:
\emph{fresh}, the initial weights of the stream run with a fresh optimizer;
\emph{early}, the end of task~0 with its optimizer state; and
\emph{late}, the end of task~49 with its optimizer state.
All starting points replace their data and dropout random streams with the same probe streams, so they
see identical batches and masks (checked by hashes). Because the probe task is the same for all starting
points, task difficulty cancels in the late-versus-early comparison. The primary metric is
\emph{probe AUC}, the mean held-out probe accuracy over the 20 evaluation points. We also report the
first-50-step mean, final accuracy, and the accuracy drop on the starting point's last training task.

\paragraph{Interventions and controls.}
Each branch restores the late checkpoint and applies one intervention to the optimizer state.
At the moment of intervention we verify bitwise that weights, BatchNorm buffers, random-number states,
stream position, hyperparameters and train- and eval-mode outputs equal those of the untouched
checkpoint. The interventions are:
\keep;
\texttt{reset\_m} ($m{\leftarrow}0$);
\texttt{reset\_v} ($v{\leftarrow}0$);
\texttt{reset\_both} (both moments);
\texttt{reset\_t} (counter only, as in \citealp{ellis2024adam}); and
two \emph{mis-corrected} artifact controls, \texttt{reset\_v} and \texttt{reset\_both} with the counter
left unchanged (``shared-keep'').
The first three resets use per-moment counters (Section~\ref{sec:analysis}). For every primary
intervention $X$, the \emph{update-norm control} \texttt{keepdir\_$X$mag} keeps the untouched optimizer
state and computes its own update. At each step it rescales that update, tensor by tensor, to the norm
that $X$'s branch applied at the same step. The control therefore needs a completed run of $X$ first
(one extra 200-step probe) and reads privileged per-step norms. It is a diagnostic, not a deployable
method. All interventions are applied at a pre-specified oracle boundary, the end of the stream. This is
not a boundary-free or online method.

\paragraph{Decision rules (fixed before running).}
The minimum interesting effect is 0.02 probe AUC.
The phenomenon is \emph{established} for a condition if the mean late-minus-early AUC is at most $-0.02$
and every seed is negative.
An intervention $X$ has an \emph{effect} if its mean difference to \keep{} is at least $0.02$ and every
seed is positive. A \emph{difference remains} if, in addition, $|X-\texttt{keepdir\_}X\texttt{mag}|$
is at least $0.02$ on average with a consistent sign across seeds.
We use three test seeds (3000--3002). These seeds were already inspected in an earlier smoke run, so
this is an exploratory study, not a confirmatory one. The run used \RunCpuSeconds{}\,s of CPU under a
\RunCpuCap{}\,s cap and passed all \RunCellsTotal{} planned cells (code \texttt{\RunCommit}).
